\newcommand{\ipPlaceHolder}{
    \foreach \i in {1,...,3}{\TextField[name=ip\i, width=1.5cm,print,borderwidth=0, backgroundcolor={0.9 0.9 0.9}]{}.}\TextField[name=ip4, width=1.5cm,print,borderwidth=0, backgroundcolor={0.9 0.9 0.9}]{}
}
\chapter{Anwendung: Chat-Kommunikation über einen Relay-Server}

In diesem Kapitel werden wir die Kenntnisse des vorausgehenden Kapitels anwenden, um ein einfaches Chat-Programm zu entwickeln, das über einen Relay-Server kommuniziert. Ziel ist es, dass zwei oder mehr Clients Nachrichten über das Internet austauschen können.


% \chapter{Lernaufgabe: Netzwerke}
% \textbf{Ziel:} Am Ende soll ein funktionsfähiges Programm stehen, mit dem die Lernenden in Echtzeit Nachrichten austauschen können.

% \textbf{Problembezug:} \enquote{Wie kann ich mit einer anderen Person eine Nachricht über das Internet austauschen?}

% \begin{mainBox}[title=Komplexe Lernaufgabe: Chat-Anwendung]
%     Die \textbf{Haupt-Lernaufgabe} lautet: Entwickeln Sie in Kleingruppen eine \textbf{funktionierende Chat-Anwendung} in Python, die über einen Relay-Server läuft. Am Ende sollen die Gruppen miteinander in Echtzeit Nachrichten austauschen können. Diese Haupt-Lernaufgabe ist in mehrere \textbf{Learning Tasks} unterteilt, die aufeinander aufbauen und jeweils durch unterstützende Informationen, Just-in-time Hilfen und Part-task Practice (Teilübungen) begleitet werden.
% \end{mainBox}

% \begin{scaffoldBox}[title=Scaffolding der Lernaufgaben]
%     Die Unterrichtssequenz besteht aus drei aufeinander aufbauenden Learning Tasks, welche jeweils durch vorausgehende, unterstützende Informationen, Just-in-time Hilfen und Part-task Practice (Teilübungen) begleitet werden. Dieses \textit{Scaffolding} jedes Learning Tasks wird jeweils in einer blauen Box hervorgehoben, ebenso wie das fortschreitende Zurückweichen (\textit{Fading}) der Unterstützung.
% \end{scaffoldBox}

% \begin{table}[H]
%     \centering
%     \small
%     \renewcommand{\arraystretch}{1.2}
%     \begin{tabular}{p{0.33\linewidth}|p{0.33\linewidth}|p{0.33\linewidth}}
%         \textbf{Unterstützende Informationen}                                                  & \textbf{Just-in-time Informationen} & \textbf{Part-task Practice (Teilübungen)} \\
%         \midrule

%         \begin{itemize}
%             \item Visualisierung der Client-Server-Architektur (Skript, Slides)
%             \item Einführung in grundlegende Begriffe: IP-Adresse, Port, Socket
%             \item Analogie: \enquote{Postsystem} (Adresse $\rightarrow$ Empfänger)
%             \item Infoblätter / Moodle-Materialien mit Hintergrundwissen
%         \end{itemize}                 &
%         \begin{itemize}
%             \item Code-Snippets für \lstinline|socket.connect|, \lstinline|send|, \lstinline|recv|
%             \item Anleitung zur Fehlerbehandlung bei Verbindungsabbrüchen
%             \item Debugging-Tipps, wenn Nachrichten nicht ankommen
%             \item Schritt-für-Schritt-Hilfekarten (nur bei Bedarf zugänglich)
%         \end{itemize} &
%         \begin{itemize}
%             \item \textbf{Erste Übung:} Zwei Python-Skripte auf demselben Rechner verbinden
%             \item Übung: Eine \enquote{Hallo}-Nachricht an den Server senden und Antwort empfangen
%             \item Fehlerübungen: Verbindungsfehler identifizieren und beheben
%             \item Mini-Tests zu IP/Port-Adressen und deren Bedeutung
%         \end{itemize}
%     \end{tabular}
%     \caption{Didaktische Struktur der Lernaufgabe: Chat-Kommunikation}
% \end{table}

% \section{Learning Task 1: Lokale Verbindung via \texorpdfstring{\texttt{localhost}}{localhost}}
\section{Lokale Chat-Kommunikation via \texorpdfstring{\texttt{localhost}}{localhost}}

In dieser ersten Aufgabe machen wir uns mit Netzwerk-Programmierung in Python vertraut, indem wir eine einfache Chat-Kommunikation auf demselben Rechner implementieren. Wir verwenden dazu die Loopback-Adresse \texttt{localhost} (IP: \texttt{127.0.0.1}), welche auf den eigenen Rechner verweist. Dies ermöglicht es uns, die Grundlagen der Socket-Programmierung zu erlernen, ohne uns um Netzwerkverbindungen zwischen verschiedenen Geräten kümmern zu müssen.

Ein \textbf{Socket} ist ein Endpunkt für die Kommunikation zwischen zwei Computern über ein Netzwerk, in welchem folgende Daten definiert werden müssen:
\begin{itemize}
    \item \textbf{IP-Adresse}: Die Adresse des Rechners, mit dem kommuniziert werden soll. Für lokale Tests verwenden wir \texttt{localhost} oder \texttt{127.0.0.1}.
    \item \textbf{Port}: Eine Nummer, die einen spezifischen Dienst oder eine Anwendung auf dem Rechner identifiziert. Ports reichen von 0 bis 65535, wobei Ports unter 1024 oft für Systemdienste reserviert sind.
    \item \textbf{Protokoll}: Das Kommunikationsprotokoll, das verwendet wird, z.B. \gls{TCP} (Transmission Control Protocol) für zuverlässige Verbindungen oder \gls{UDP} für schnellere, aber unzuverlässige Verbindungen.
\end{itemize}

% \begin{scaffoldBox}
%     \textbf{Vorausgehende Konzepte}: Die SuS kennen bereits die historischen Hintergründe von Netzwerken, die Grundkonzepte von Netzwerkschichten, IP-Adressen, Sequenzierung und Ports, sowie die Client-Server-Architektur aufgrund einer vorangehenden Unterrichtssequenz. In der Vorausgehenden Unterrichtssequenz haben sich die SuS mit den grundlegenden Funktionen und der Bedeutung von Netzwerken auseinandergesetzt, indem Sie ein Spiel gespielt haben, bei dem sie Nachrichten über ein simuliertes Netzwerk via ausgedruckte Kärtchen zusenden mussten (s. \cref{chapter:netzwerkspiel}). Dabei haben sie die Rolle von IP-Adressen und Ports kennengelernt und verstanden, wie Datenpakete in einem Netzwerk weitergeleitet werden.

%     \textbf{Neue Konzepte}: Client-Server-Architektur, IP-Adresse, Port, Socket, localhost.

%     Die neuen Begriffe werden zuerst anhand eines Videos (\href{https://www.youtube.com/watch?v=PBWhzz_Gn10}{Warriors of the Net}) und einer Analogie (Postsystem) eingeführt. Anschliessend wird die Grundstruktur eines Socket-Programms in Python vorgestellt. Die SuS lösen dabei einige vorbereitete Teilübungen (Part-task Practice), um die Konzepte zu verinnerlichen, und lösen danach die komplexe Lernaufgabe, bei der eine lokale Chat-Kommunikation über zwei Python-Skripte realisiert wird.
% \end{scaffoldBox}

\begin{myexercise}[Chatten auf demselben Rechner]\label{ex:chat-localhost}

    \textbf{Aufgabe:} Schreiben Sie zwei Python-Programme, welche auf demselben Rechner laufen: \texttt{server\_localhost.py} und \texttt{client\_localhost.py}. Der Client verbindet sich mit dem Server und schickt eine kurze Nachricht: \enquote{Hallo Server, hier ist der Client}. Der Server empfängt die Nachricht und gibt sie auf der Konsole aus, zudem antwortet er an den Client \enquote{Hallo Client, hier ist der Server}.

    Folgendes Video führt Sie durch alle Schritte: \href{https://youtu.be/1ECIEdc3TdA}{Link}. Alternativ finden Sie unten eine Schritt-für-Schritt-Anleitung.

    \textbf{Verwenden Sie folgende Grundstruktur sowohl in den Dateien \texttt{server\_localhost.py} als auch \texttt{client\_localhost.py}:}

    \begin{lstlisting}
import socket # sockets-Library importieren, damit wir Sockets erstellen können

PORT = ... # irgendein Port für den Server / Client definieren (z.B. 12345)
IP_ADDRESS = ... # lokale IP-Adresse des eigenen Computers ("localhost" oder "127.0.0.1")
    \end{lstlisting}

    Danach können Sie den Server-Socket und Client-Socket erstellen. Ein Socket wird mit \lstinline|socket.socket()| erstellt. Die Argumente geben den Adresstyp und das Protokoll an. In unserem Fall bedeutet \lstinline|socket.AF_INET|, dass IPv4-Adressen verwendet werden, und \lstinline|socket.SOCK_STREAM| steht für eine \gls{TCP}-Verbindung.

    \begin{lstlisting}
# Erstelle einen TCP-IP Socket
server = socket.socket(socket.AF_INET, socket.SOCK_STREAM)
    \end{lstlisting}

    Danach muss der Server-Socket an eine IP-Adresse und einen Port gebunden werden, damit er auf eingehende Verbindungen lauschen kann. Dies geschieht mit der Methode \lstinline|bind()|.
    \begin{lstlisting}
# Binde den Socket an localhost und den definierten Port
server.bind((IP_ADDRESS, PORT))
    \end{lstlisting}

    Der Server muss dann auf eingehende Verbindungen warten, was mit der Methode \lstinline|listen()| geschieht. Danach kann eine Verbindung mit \lstinline|accept()| akzeptiert werden.
    \begin{lstlisting}
server.listen(1)  # Warte auf eingehende Verbindungen (maximal 1 in der Warteschlange)
conn, addr = server.accept()  # Akzeptiere eine eingehende Verbindung
    \end{lstlisting}

    Der Client verbindet sich mit dem Server mittels der Methode \lstinline|connect()|:
    \begin{lstlisting}
# Verbinde den Socket mit dem Server auf localhost und dem definierten Port
client.connect((IP_ADDRESS, PORT))
    \end{lstlisting}
    Nun können Sie im Server und im Client Nachrichten senden und empfangen, indem Sie die Methoden \lstinline|send()| und \lstinline|recv()| verwenden. Zuerst müssen Sie jedoch Ihre Nachricht als String definieren:
    \begin{lstlisting}
# Beispiel-Nachricht
nachricht = "Hallo Server, hier ist der Client"
    \end{lstlisting}
    Danach kann diese Nachricht an den Server gesendet werden (verwenden Sie \texttt{.encode()}, um den String in Bytes umzuwandeln):
    \begin{lstlisting}
client.send(nachricht.encode())  # Sende die Nachricht an den Server
    \end{lstlisting}
    Um eine Nachricht zu empfangen, verwenden Sie die Methode \lstinline|recv()|, welche die Anzahl der Bytes angibt, die empfangen werden sollen (z.B. 1024 Bytes):
    \begin{lstlisting}
empfangene_nachricht = conn.recv(1024).decode()  # Empfange bis zu 1024 Bytes
print(f"Empfangen: {empfangene_nachricht}")  # Zeige die empfangene Nachricht an
    \end{lstlisting}
    \textbf{Hinweis:} Vergessen Sie nicht, die Sockets am Ende mit \lstinline|close()| zu schliessen, um Ressourcen freizugeben:
    \begin{lstlisting}
conn.close()
server.close()
client.close()
     \end{lstlisting}

    Lassen Sie \texttt{server\_localhost.py} auf der Konsole laufen und führen Sie dann \texttt{client\_localhost.py} aus, indem Sie in VS Code auf den Dropdown-Button rechts vom Run-Button clicken und \enquote{Run Python File in Dedicated Terminal} anklicken.
    \begin{figure}[H]
        \begin{tikzpicture}[
                base/.style={rectangle, draw, minimum width=2.5cm, minimum height=1cm, text=white, rounded corners, text width=2.8cm, align=center, inner sep=8pt},
                server/.style={base, fill=RoyalBlue},
                client/.style={base, fill=ForestGreen},
                arrow/.style={-Stealth, thick, rounded corners},
                code/.style={font=\ttfamily\small},
                socket/.style={circle, draw, fill=yellow!10, minimum size=0.5cm}
            ]

            % Socket symbol inside server rectangle
            \node[client] (client) at (0, 0) {Client\\\texttt{198.168.0.10}};
            \node[server] (server) at (10, 0) {Server\\\texttt{198.168.0.11}};
            \node[socket] (server_socket) at (server.north west) {\large \emoji{pig-nose}};

            % Arrows with labels

            \node[rectangle, draw, fill=yellow!5, above right=.2cm and .6cm of server_socket, minimum width=2.2cm, minimum height=1.2cm, font=\small, align=center] {
                \textbf{Socket:}\\
                IP: \texttt{198.168.0.11}\\
                Port: \texttt{12345}\\
                Protocol: \texttt{TCP}
            };
            \draw[arrow, loop, min distance=2cm, in=120, out=60] (server_socket.north) to node[above, yshift=10pt, code, align=center, text width=2cm] {bind()\\ listen()\\ accept()} (server_socket.north);
            \node[socket] (client_socket) at (client.north east) {};
            \draw[arrow] (client_socket) to[bend left=40] node[above, yshift=3pt, code] {connect() \faPlug{}} (server_socket);
            \draw[arrow] (client_socket) -- node[below, code] {send("Hallo Server...")} (server_socket);
            \draw[arrow, dashed] (server_socket) -- node[above, yshift=3pt, code, xshift=10pt] {recv()} (client_socket);
            \draw[arrow, dashed, bend left=40] (server_socket.south) to node[below, yshift=-8pt, code, xshift=-10pt] {send("Hallo Client...")} (client_socket.south);
            % State boxes
            \node[rectangle, draw, fill=gray!10, below=3cm of client.east, text width=7cm, align=left, font=\small] (client_state) {
                \textbf{Client:} \\
                1. Socket erstellen (\texttt{.socket()})\\
                2. Mit Server-Socket verbinden (\texttt{.connect()}) \\
                3. Nachricht senden (\texttt{.send()})\\
                4. Antwort empfangen (\texttt{.recv()})\\
            };

            \node[rectangle, draw, fill=gray!10, below=3cm of server.east, xshift=-2cm, text width=7cm, align=left, font=\small] (server_state) {
                \textbf{Server:} \\
                1. Socket erstellen (\texttt{.socket(), .bind()})\\
                2. Warten auf Verbindungen (\texttt{.listen()}) \\
                3. Accept Verbindung (\texttt{.accept()}) \\
                4. Nachricht empfangen (\texttt{.recv()}) \\
                5. Antwort senden (\texttt{.send()})
            };
        \end{tikzpicture}
        \centering
        \caption{Ablauf einer Socket-Kommunikation zwischen Client und Server auf \texttt{localhost}}
    \end{figure}
    Ein Host kann auch gleichzeitig einen Server-Dienst und einen Client-Dienst ausführen, um mit sich selber zu kommunizieren -- denn letztlich ist ein Server nur ein Programm, das auf Anfragen von Clients wartet. In der folgenden Abbildung sehen Sie einen Client, der zwei Sockets verwendet, um bidirektional mit sich selber zu kommunizieren:
    \begin{figure}[H]
        \begin{tikzpicture}[
                base/.style={rectangle, draw, minimum width=2.5cm, minimum height=1cm, text=white, rounded corners, text width=2.8cm, align=center, inner sep=8pt},
                client/.style={base, fill=ForestGreen},
                socket/.style={circle, draw, fill=yellow!10, minimum size=0.6cm},
                arrow/.style={-Stealth, thick, rounded corners}
            ]

            \node[client] (client) at (5, 0) {Client\\\texttt{localhost:54321}};
            \node[socket] (socket1) at (3.5, 0) {};
            \node[socket] (socket2) at (6.5, 0) {};
            \draw[arrow] (socket1) to[bend left=60] (socket2);
            \draw[arrow] (socket2) to[bend left=60] (socket1);
            \node[rectangle, draw, fill=yellow!5, above right=.2cm and .6cm of socket1, minimum width=2.2cm, minimum height=1.2cm, font=\small, align=center] {
                \textbf{Socket:}\\
                IP: \texttt{localhost}\\
                Port: \texttt{54321}\\
                Protocol: \texttt{TCP}
            };
        \end{tikzpicture}
        \centering
        \caption{Client mit zwei Sockets für bidirektionale Kommunikation}
    \end{figure}
\end{myexercise}

\begin{myexercise}
    Erweitern Sie nun Ihr Programm, indem der Client eine Nachricht schickt, und der Server antwortet mit \enquote{Nachricht erhalten!}. Sie können hierfür serverseitig den Befehl \lstinline|conn.send()| verwenden, um die Antwort zu senden, und clientseitig \lstinline|client.recv()|, um die Antwort im Client zu empfangen und anzuzeigen.
\end{myexercise}

\begin{myattention}[Gebundene Ports]
    Wenn Sie den Server beenden, wird der Port unter Umständen für kurze Zeit weiterhin als \enquote{belegt} angezeigt. Dies liegt daran, dass das Betriebssystem den Port für eine gewisse Zeit reserviert, um sicherzustellen, dass alle Daten korrekt übertragen wurden. Wenn Sie versuchen, den Server sofort neu zu starten und denselben Port zu verwenden, kann es zu einem Fehler kommen, da der Port noch als belegt gilt. In diesem Fall können Sie folgende Zeile hinzufügen, um den Port sofort wiederverwendbar zu machen:
    \begin{lstlisting}
server.setsockopt(socket.SOL_SOCKET, socket.SO_REUSEADDR, 1)
    \end{lstlisting}
    Fügen Sie diese Zeile direkt nach der Erstellung des Server-Sockets ein, also nach folgender Zeile:
    \begin{lstlisting}
server = socket.socket(socket.AF_INET, socket.SOCK_STREAM)
    \end{lstlisting}
\end{myattention}

\begin{myexercise}[Verständnisfragen]
    \begin{itemize}
        \item Was geschieht, wenn in \texttt{server\_localhost.py} die IP-Adresse \lstinline|"localhost"| durch \lstinline|"127.0.0.1"| ersetzt wird?
        \item Was geschieht, wenn eine andere Port-Nummer als im Client und Server verwendet wird? Funktioniert die Verbindung? Weshalb (nicht)?
        \item Warum muss der Server gestartet werden, bevor der Client versucht, sich zu verbinden?
        \item Was passiert, wenn der Server eine Nachricht sendet, aber der Client keine Antwort erwartet?
    \end{itemize}
    \begin{myanswer}
        \begin{itemize}
            \item \lstinline|"localhost"| und \lstinline|"127.0.0.1"| sind äquivalent und verweisen beide auf die lokale IP-Adresse des eigenen Rechners.
            \item Die Verbindung funktioniert nicht, da der Client versucht, sich mit einem Port zu verbinden, auf dem der Server nicht lauscht. Beide müssen denselben Port verwenden.
            \item Der Server muss zuerst gestartet werden, damit er auf Verbindungsanfragen lauschen kann. Wenn der Client zuerst startet, gibt es keinen Prozess, der die Verbindung annimmt, was zu einem Fehler führt.
            \item Wenn der Server eine Nachricht sendet, aber der Client keine Antwort erwartet, kann es sein, dass die Nachricht verloren geht oder der Client sie nicht verarbeitet, was zu Kommunikationsproblemen führen kann.
        \end{itemize}
    \end{myanswer}
\end{myexercise}


\begin{myexercise}[Reflexionsfragen]
    Notieren Sie nach Abschluss aller vorausgehenden Aufgaben kurze Antworten in Ihr Lernjournal. Tauschen Sie sich danach mit einer Person aus und ergänzen Sie fehlende Punkte.

    \textbf{Fragen:}
    \begin{enumerate}
        \item Erklären Sie in \emph{eigenen Worten}: Was ist ein Socket? Welche zwei Operationen haben Sie verwendet?
        \item Unterschied Socket vs. Port?
        \item Unterschied \gls{IP}-Adresse vs. Port (verwenden Sie die Post-Analogie).
        \item Warum muss der Server vor dem Client gestartet werden? Welche konkrete Fehlermeldung erscheint sonst?
        \item Weshalb muss \lstinline|bind()| nur einmal aufgerufen werden?
        \item Skizzieren Sie den Ablauf Ihrer Nachricht (Methodenaufrufe in Reihenfolge) vom Client zum Server.
        \item Nennen Sie mindestens einen Debugging-Schritt, falls keine Antwort zurückkommt.
        \item Formulieren Sie eine mögliche Ursache für \texttt{ConnectionRefusedError}.
        \item Kurzer Selbstcheck (Ja/Nein):
              \begin{itemize}
                  \item Ich kann den Code ohne Vorlage erneut schreiben.
                  \item Ich verstehe jede Zeile von \texttt{server.py}.
                  \item Ich weiss, wann \texttt{recv} blockiert.
              \end{itemize}
        \item Formulieren Sie ein persönliches Lernziel für Task 2 (ein Satz).
    \end{enumerate}

    \begin{myanswer}
        \begin{itemize}
            \item Socket: Endpunkt einer TCP-Verbindung; genutzt: \texttt{socket()}, \texttt{bind()}, \texttt{listen()}, \texttt{accept()}, \texttt{connect()}, \texttt{send()}, \texttt{recv()}, \texttt{close()}.
            \item Socket = Verbindungspunkt (Python-Objekt); Port = Nummer zur Identifikation eines Dienstes auf einem Rechner.
            \item \gls{IP} = Adresse des Rechners; Port = „Tür"/Dienst auf diesem Rechner (Post: Hausnummer vs. Zimmernummer).
            \item Server zuerst: Sonst kein Prozess lauscht $\rightarrow$ \texttt{ConnectionRefusedError: [Errno 61] Connection refused}.
            \item \texttt{bind()} weist Socket eine Adresse/Port zu; nur einmal nötig, da Socket danach gebunden ist.
            \item Ablauf:
                  \begin{itemize}
                      \item \textbf{Client}: \texttt{socket()} $\rightarrow$ \texttt{connect()} $\rightarrow$ \texttt{send()};
                      \item \textbf{Server}: \texttt{socket()} $\rightarrow$ \texttt{bind()} $\rightarrow$ \texttt{listen()} $\rightarrow$ \texttt{accept()} $\rightarrow$ \texttt{recv()} $\rightarrow$ (optional \texttt{send()} Antwort).
                  \end{itemize}

            \item Debugging: Prüfen Port, laufender Server, Firewall, Print-Statements vor/nach \texttt{send/recv}, \texttt{netstat} zur Belegung.
            \item Ursache \texttt{ConnectionRefusedError}: Server nicht gestartet / falscher Port / Firewall blockiert.
            \item Blockierendes \texttt{recv}: Wartet bis Daten eintreffen oder Gegenstelle schliesst.
            \item Lernziel-Beispiel: „Ich möchte in Task 2 mehrere Nachrichten nacheinander austauschen können."
        \end{itemize}
    \end{myanswer}
\end{myexercise}

% \section{Learning Task 2: Chat-Kommunikation über ein lokales Netzwerk}
\section{Chat-Kommunikation über ein lokales Netzwerk (LAN)}

% \begin{scaffoldBox}
%     \textbf{Neue Konzepte}: IP-Adresse im lokalen Netzwerk, Firewall-Einstellungen.

%     Die SuS lernen, wie sie die eigene IP-Adresse im lokalen Netzwerk ermitteln können (z. B. über den Befehl \texttt{ipconfig} in der Windows-Eingabeaufforderung oder \texttt{ifconfig} im Terminal auf macOS/Linux). Zudem wird erklärt, wie Firewall-Einstellungen den Datenverkehr beeinflussen können. Die Teilübungen (Part-task Practice) helfen den SuS, ihre Programme anzupassen und Verbindungsprobleme zu lösen. Am Ende steht die komplexe Lernaufgabe, bei der die Chat-Kommunikation über zwei Geräte im selben Netzwerk realisiert wird.

%     \textbf{Technische Herausforderungen}: Die Lehrperson sollte vorab testen, ob die Firewall-Einstellungen auf den Schulrechnern die Kommunikation über Sockets erlauben. Gegebenenfalls müssen Ausnahmen für die verwendeten Ports eingerichtet werden oder ein eigenes Netzwerk (z. B. über einen mobilen Hotspot oder einen eigenen WLAN-fähigen Router) genutzt werden.
% \end{scaffoldBox}

\begin{myexercise}[Meine private und meine öffentliche \gls{IP}-Adresse herausfinden]
    Finden Sie die lokale \gls{IP}-Adresse Ihres Laptops heraus. Dies ist die IP, welche Sie innerhalb des lokalen Netzwerks (z.B. Schul-WLAN) verwenden. Verwenden Sie dazu folgende Anleitung:
    \begin{itemize}
        \item[\faApple] \textbf{MacOS}: Öffnen Sie das \textbf{Terminal} (dieses können Sie via Spotlight-Suche finden: \keys{\cmd + \SPACE}). Tippen Sie im Terminal: \commandbox*|ipconfig getifaddr en0|
		\item[\faWindows] \textbf{Windows}: Öffnen Sie die \textbf{Eingabeaufforderung} (diese können Sie finden durch Drücken der \keys{\faWindows}-Taste, danach Suche nach dem Programm \texttt{cmd} oder Eingabeaufforderung). Tippen Sie in der Eingabeaufforderung: \commandbox*|ipconfig|
    \end{itemize}

    Ihre private \gls{IP}-Adresse sollte in etwa so aussehen: \texttt{192.168.x.x} oder \texttt{10.x.x.x}.

    Meine private \gls{IP}-Adresse ist: \ipPlaceHolder

    Notieren Sie sich auch noch ihre öffentliche \gls{IP}-Adresse, welche Sie im Internet verwenden. Diese finden Sie z.B. auf der Webseite \url{https://www.whatismyip.com/}.

    Meine öffentliche \gls{IP}-Adresse ist: \ipPlaceHolder
\end{myexercise}

\begin{myexercise}
    Vergleichen Sie die privaten und die öffentlichen \gls{IP}-Adressen mit denjenigen Ihrer Pultnachbarin. Welche Unterschiede und Gemeinsamkeiten gibt es? Was könnte der Grund dafür sein? Würden Sie für die gemeinsame Chat-Kommunikation innerhalb des gleichen Netzwerks die private oder die öffentliche \gls{IP}-Adresse verwenden?

    \begin{myanswer}
        Die privaten \gls{IP}-Adressen der beiden Rechner sind wahrscheinlich ähnlich aber nicht gleich (z.B. beide beginnen mit \texttt{192.168.x.x}), da sie sich im selben lokalen Netzwerk befinden. Die öffentlichen \gls{IP}-Adressen sind hingegen identisch, da beide Rechner über denselben Internetanschluss kommunizieren. Für die Chat-Kommunikation innerhalb des gleichen Netzwerks sollte die private \gls{IP}-Adresse verwendet werden, da sie direkt auf den jeweiligen Rechner im LAN verweist, während die öffentliche \gls{IP}-Adresse nur für die Kommunikation über das Internet relevant ist.
    \end{myanswer}
\end{myexercise}

\begin{myattention}[Firewall-Einstellungen]
    In einigen Fällen können Firewall-Einstellungen die Kommunikation über Sockets blockieren, insbesondere wenn Sie versuchen, Verbindungen von einem anderen Gerät im Netzwerk herzustellen. Wenn Sie Verbindungsprobleme haben, überprüfen Sie die Firewall-Einstellungen auf beiden Geräten und stellen Sie sicher, dass der Port, den Sie verwenden, für eingehende Verbindungen geöffnet ist. In manchen Fällen müssen Sie eine Ausnahme für Ihr Python-Programm oder den verwendeten Port hinzufügen.

    Für Windows: Gehen Sie zu \textbf{Einstellungen} → \textbf{Netzwerk & Internet} → \textbf{Eigenschaften} und stellen Sie sicher, dass Ihr Netzwerkprofil auf \textbf{Privat} eingestellt ist, damit die Firewall Verbindungen von anderen Geräten im selben Netzwerk zulässt.
\end{myattention}

\begin{myexercise}[Chatten im lokalen Netzwerk]\label{ex:chat-localnet}

    \textbf{Aufgabe:} Erweitern Sie Ihre Codes aus \cref{ex:chat-localhost}, so dass ihre Programme auf \textbf{verschiedenen Rechnern im selben lokalen Netzwerk} laufen. Der Server lauscht auf seiner eigenen, privaten \gls{IP}-Adresse und einem Port (z.B. 12345). Der Client verbindet sich mit der \gls{IP}-Adresse des Servers und schickt eine kurze Nachricht (\enquote{Hallo Server}). Der Server empfängt die Nachricht und gibt sie auf der Konsole aus.

    \textbf{Schritte:}
    \begin{enumerate}
        \item Ermitteln Sie die private \gls{IP}-Adresse ihres Rechners und desjenigen ihrer Pultnachbarin (s. \cref{ex:meine-ip}). Bestimmen Sie, wer von Ihnen die Rolle des Servers und wer die Rolle des Clients übernimmt.
        \item Passen Sie die \texttt{bind}- und \texttt{connect}-Befehle in beiden Skripts entsprechend an.
        \item Starten Sie \texttt{server.py} auf Rechner A und \texttt{client.py} auf Rechner B und tauschen Sie eine Nachricht aus.
    \end{enumerate}
\end{myexercise}

\begin{myexercise}[Endlos-Chat]\label{ex:endlos-chat}
    Im echten Leben wollen wir nicht nur eine einzige Nachricht austauschen, sondern mehrere Nachrichten hin- und herschicken. Erweitern Sie Ihre Programme so, dass der Client und der Server in einer Endlosschleife Nachrichten austauschen können, bis einer der beiden die Verbindung mit \keys{\cmd + C} (MacOS) bzw. \keys{\ctrl + C} (Windows) abbricht. Verwenden Sie dafür folgende Grundstruktur im Server:
    \begin{lstlisting}
try:
    # Endlosschleife für den Nachrichtenaustausch
    while True:
        # Nachricht empfangen
        nachricht = conn.recv(1024).decode()
        print("Empfangen:", nachricht)

        # Nachricht senden
        antwort = input("Antwort an den Client: ")
        conn.send(antwort.encode())
except KeyboardInterrupt:
    print("Verbindung wird geschlossen.")
    # Socket schliessen
    conn.close()
    server.close()
\end{lstlisting}

Passen Sie den Client entsprechend an, damit auch dieser in einer Endlosschleife Nachrichten senden und empfangen kann. Testen Sie die Kommunikation, indem Sie abwechselnd Nachrichten eingeben und empfangen.
\end{myexercise}

\begin{myexercise}[Reflexionsfragen]
    Notieren Sie nach Abschluss von Learning Task 2 (max. 5 Minuten) kurze Antworten in Ihr Lernjournal. Tauschen Sie sich danach mit einer Person aus und ergänzen Sie fehlende Punkte.

    \textbf{Fragen:}
    \begin{enumerate}
        \item Wie unterscheidet sich die \gls{IP}-Adresse im lokalen Netzwerk von der \texttt{localhost}-Adresse?
        \item Warum ist es wichtig, die richtige \gls{IP}-Adresse des Servers in der Datei \texttt{client.py} zu verwenden?
        \item Welche Fehlermeldung erscheint, wenn der Client versucht, sich mit einer falschen \gls{IP}-Adresse zu verbinden?
        \item Kurzer Selbstcheck (Ja/Nein):
              \begin{itemize}
                  \item Ich verstehe jede Zeile von \texttt{server.py}.
                  \item Ich weiss, wie ich die lokale IP-Adresse meines Rechners finde.
              \end{itemize}
        \item Ein persönliches Lernziel für Task 3 (ein Satz).
        \item Welche Herausforderungen könnten auftreten, wenn Sie versuchen, über das Internet zu kommunizieren, anstatt nur im lokalen Netzwerk?
    \end{enumerate}
    \begin{myanswer}
        \begin{enumerate}
            \item Unterschied: \texttt{localhost} (127.0.0.1) verweist nur auf den eigenen Rechner; die lokale IP (z.B. 192.168.x.x) identifiziert den Rechner im LAN für andere Geräte.
            \item Richtige \gls{IP} nötig, sonst versucht der Client einen nicht existierenden Host oder anderen Rechner zu erreichen $\rightarrow$ keine oder falsche Verbindung.
            \item Mögliche Fehlermeldungen: \texttt{ConnectionRefusedError} (Dienst nicht aktiv), Timeout (keine Antwort), \texttt{OSError: No route to host}.
            \item Selbstcheck: Code reproduzierbar? Jede Zeile verstanden? IP finden gelingt? Wenn Nein $\rightarrow$ gezielt nacharbeiten.
            \item Lernziel Task 3: \enquote{Ich kann ein Client-Server-Programm schreiben, um gezielt Nachrichten an einen benannten Server zu senden.}
            \item Herausforderungen Internet: NAT und Port-Forwarding nötig; dynamische öffentliche IP; zusätzliche Latenz; Firewalls/Provider-Restriktionen; Sicherheit (Verschlüsselung, Authentifizierung); DNS-Auflösung.
        \end{enumerate}
    \end{myanswer}

\end{myexercise}
\section{Chat-Kommunikation über das LAN, mittels Relay-Server}
% \section{Learning Task 3: Chat-Kommunikation über das LAN, mittels Relay-Server}
% \begin{scaffoldBox}
%     \textbf{Neue Konzepte}: Relay-Server, Client-Identifikation.

%     Die SuS lernen das Konzept eines Relay-Servers kennen, der als Vermittler zwischen mehreren Clients fungiert. Der Relay-Server empfängt Nachrichten von einem Client und leitet sie an den entsprechenden Empfänger-Client weiter. Die Teilübungen (Part-task Practice) helfen den SuS, ihre Programme anzupassen und die Logik für die Client-Identifikation zu implementieren. Am Ende steht die komplexe Lernaufgabe, bei der die Chat-Kommunikation über einen Relay-Server im lokalen Netzwerk realisiert wird.
% \end{scaffoldBox}



\subsection{Überblick: Wie funktioniert ein Chat-Client?}

Ein Chat-Client verbindet sich über das Internet mit einem Server, also einem Computer, der Nachrichten von Clients hin- und herschickt. Nachrichten werden als Text über ein das \gls{TCP}-Protokoll verschickt. Der Server kann die Nachricht an andere Clients weiterleiten oder – wie beim Echo-Server – einfach zurückschicken.

\begin{figure}[H]
    \begin{tikzpicture}[
            base/.style={rectangle, draw, minimum width=2cm, minimum height=1cm, text=white, rounded corners, text width=2.4cm, align=center,inner sep=8pt},
            server/.style={base, fill=RoyalBlue},
            client/.style={base, fill=ForestGreen},
            arrow/.style={-Stealth, thick, rounded corners, text width=4cm, align=center}
        ]

        % Nodes
        \node[client] (client1) {Client 1\\\texttt{192.168.0.20}};
        \node[server, right=6cm of client1] (server) {Echo-Server\\\texttt{192.168.0.30}};

        % Arrows
        \draw[arrow] (client1) -- node[above, yshift=4pt] {\texttt{192.168.0.20:Hallo Max!}} (server);
        \draw[arrow] (server) -- node[below, yshift=-4pt] {\texttt{192.168.0.20:Hallo Max!}} (client1);

    \end{tikzpicture}
    \centering
    \caption{Ablauf einer Chat-Nachricht über den Echo-Server}
\end{figure}

\begin{figure}[H]
    \begin{tikzpicture}[
            base/.style={rectangle, draw, minimum width=2cm, minimum height=1cm, text=white, rounded corners, text width=2.4cm, align=center,inner sep=8pt},
            server/.style={base, fill=RoyalBlue},
            client/.style={base, fill=ForestGreen},
            arrow/.style={-Stealth, thick, rounded corners, text width=4cm, align=center}
        ]

        % Nodes
        \node[client] (client1) {Client 1\\\texttt{192.168.0.20}};
        \node[server, below right=of client1] (relay) {Relay-Server\\\texttt{192.168.0.40}};
        \node[client, above right=of relay] (client2) {Client 2\\\texttt{192.168.0.21}};

        % Arrows
        \draw[arrow] (client1) -- node[midway, left] {\texttt{192.168.0.21:Hallo!}} (relay);
        \draw[arrow] (relay) -- node[midway, right] {\texttt{192.168.0.21:Hallo!}} (client2);

    \end{tikzpicture}
    \centering
    \caption{Kommunikation über einen Relay-Server: Die Nachricht wird vom Client an den Relay-Server geschickt und von dort an den Ziel-Client weitergeleitet.}
\end{figure}
Im folgenden werden wir die Kommunikation über einen Relay-Server umsetzen. So können Sie Nachrichten an MitschülerInnen schicken, ohne deren \gls{IP}-Adresse zu kennen.

\subsection{Aufgaben: Schritt für Schritt zum Chat-Client}

\begin{myexercise}[Nachricht zusammensetzen]
    Ändern Sie Ihr Client-Programm so ab, dass jede Nachricht in folgendem Format verschickt wird: \texttt{IP:Nachricht}, wobei \texttt{IP} die IP-Adresse des Empfängers ist, und \texttt{Nachricht} die eigentliche Chat-Nachricht. Definieren Sie dafür drei Variablen:
    \begin{itemize}
        \item \texttt{ip}: die lokale IP-Adresse Ihres Mitschülers/Ihrer Mitschülerin als Text
        \item \texttt{message}: die Nachricht, die Sie schicken möchten
        \item \texttt{combined}: die kombinierte Chat-Nachricht im Format \texttt{ip:message}
    \end{itemize}
    Verwenden Sie dazu folgende Code-Vorlage:
    \begin{lstlisting}
ip = "20.161.75.213"
message = "Hallo Bob!"
combined = ip + ":" + message
print(combined)
\end{lstlisting}
\end{myexercise}

\begin{myexercise}[Benutzereingabe und Schleife]

    Erweitern Sie Ihr Programm so, dass die IP-Adresse und die Nachricht jeweils per \texttt{input()} eingegeben werden. Kombinieren Sie beide wie oben und geben Sie das Ergebnis aus. Bauen Sie eine Endlosschleife ein, damit Sie beliebig viele Nachrichten verschicken können.
    \begin{myanswer}
        \begin{lstlisting}[language=Python]
while True:
    ip = input("Empfänger-IP: ")
    message = input("Nachricht: ")
    combined = ip + ":" + message
    print(combined)
\end{lstlisting}
    \end{myanswer}
\end{myexercise}

\begin{myexercise}[Server-IP und Port definieren]
    Definieren Sie die IP-Adresse und den Port des Chat-Servers als Variablen \lstinline|server_ip| und \lstinline|port| (siehe Tabelle unten).
    \begin{myanswer}
        \begin{lstlisting}
server_ip = "88.198.193.239"
port = 43223
\end{lstlisting}
    \end{myanswer}
\end{myexercise}

\begin{myexercise}[Verbindung zum Server herstellen]

    Das Paket \texttt{socket} ist in Python standardmässig bereits installiert und dient dazu, um Netzwerkverbindungen herzustellen. Importieren Sie das Paket und erstellen Sie eine Funktion \lstinline|sende|, die sich mit dem Server verbindet und eine Nachricht sendet. Die Funktion soll die IP-Adresse des Servers, den Port und die Nachricht als Parameter entgegennehmen. Die Antwort des Servers soll ausgegeben werden.

    Verwenden Sie folgende Code-Vorlage: Innerhalb Ihrer Definition:
    \begin{lstlisting}[language=Python]
import socket

# IHR BISHERIGER CODE HIER

with socket.socket(socket.AF_INET, socket.SOCK_STREAM) as s: # Socket erstellen und verbinden
    s.connect((server_ip, port)) # Mit Server verbinden
    s.sendall(message.encode()) # Nachricht senden
    data = s.recv(1024) # Antwort empfangen und ausgeben
    print('Antwort vom Server:', data.decode())
\end{lstlisting}

    \begin{myanswer}
        \begin{lstlisting}[language=Python]
import socket

def send_message(server_ip, port, message):
    with socket.socket(socket.AF_INET, socket.SOCK_STREAM) as s:
        s.connect((server_ip, port))
        s.sendall(message.encode())
        data = s.recv(1024)
        print('Antwort vom Server:', data.decode())
\end{lstlisting}
    \end{myanswer}
\end{myexercise}

\begin{myexercise}[Chat-Client zusammenbauen und testen]

    Kombinieren Sie alles zu einem vollständigen Chat-Client, so dass Sie in einer Endlosschleife Nachrichten eingeben und an den Server schicken können. Die Antwort des Servers wird ausgegeben.
    \begin{myanswer}
        \begin{lstlisting}[language=Python]
import socket

server_ip = "116.203.99.55"
port = 43223

while True:
    ip = input("Empfänger-IP: ")
    message = input("Nachricht: ")
    combined = ip + ":" + message
    with socket.socket(socket.AF_INET, socket.SOCK_STREAM) as s:
        s.connect((server_ip, port))
        s.sendall(combined.encode())
        data = s.recv(1024)
        print('Antwort vom Server:', data.decode())
\end{lstlisting}
    \end{myanswer}
\end{myexercise}

\noindent
Der Client schickt eine Nachricht im Format \texttt{IP:Nachricht} an den Server. Der Echo-Server schickt exakt dieselbe Nachricht zurück. So können Sie testen, ob Ihr Client korrekt funktioniert.

% \section{Learning Task 4: Chat-Kommunikation über das Internet, mittels Relay-Server}
\section{Chat-Kommunikation über das Internet, mittels Relay-Server}

Bis jetzt haben wir nur wenige Nachrichten innerhalb des lokalen Netzwerks ausgetauscht. In der Praxis möchten wir jedoch oft mit Personen kommunizieren, die sich an einem anderen Ort befinden, also über das Internet. Zudem möchten wir einen Endlos-Fluss an Nachrichten senden und empfangen können, ohne dass die Verbindung nach jeder Nachricht getrennt wird. In diesem Kapitel lernen wir, wie dieses mit einem Relay-Server im lokalen Netzwerk sowie Python realisiert werden kann.

% \begin{scaffoldBox}
%     \textbf{Neue Konzepte}: Öffentliche IP-Adresse, NAT, Port-Forwarding, Relay-Server, DNS (Domain Name System).

%     Die SuS lernen, wie sie die öffentliche IP-Adresse ihres Netzwerks ermitteln können (z. B. über Websites wie \href{https://whatismyipaddress.com/}{whatismyipaddress.com}). Zudem wird das Konzept von NAT (Network Address Translation) und Port-Forwarding erklärt, um zu verstehen, warum direkte Verbindungen oft nicht möglich sind. Zusätzlich wird das Domain Name System (DNS) eingeführt: DNS übersetzt menschenlesbare Domain-Namen (wie \texttt{chat.example.com}) in IP-Adressen, damit Programme Server im Internet finden können. Die Teilübungen (Part-task Practice) helfen den SuS, ihre Programme anzupassen und Verbindungsprobleme zu lösen. Am Ende steht die komplexe Lernaufgabe, bei der die Chat-Kommunikation über einen Relay-Server realisiert wird.

%     \textbf{Technische Herausforderungen}: Die Lehrperson sollte vorab testen, ob die Firewall-Einstellungen auf den Schulrechnern die Kommunikation über Sockets erlauben. Gegebenenfalls müssen Ausnahmen für die verwendeten Ports eingerichtet werden oder ein eigenes Netzwerk (z. B. über einen mobilen Hotspot oder einen eigenen WLAN-fähigen Router) genutzt werden.
% \end{scaffoldBox}

\begin{myexercise}[Chatten über einen öffentlichen Relay-Server]

    \textbf{Aufgabe:} Sie verbinden sich mit einem bereits laufenden öffentlichen Relay-Server, der Nachrichten zwischen mehreren Clients weiterleitet. Ziel ist, dass Sie und Ihre Mitschüler:innen sich gegenseitig Nachrichten schicken können.

    \textbf{Schritte:}
    \begin{enumerate}
        \item Öffnen Sie \texttt{client.py} und passen Sie die \texttt{connect}-Zeile an:
              \begin{lstlisting}[language=Python]
client.connect(("88.198.193.239", 12345))
\end{lstlisting}
        \item Senden Sie eine Nachricht (z.B. \enquote{Hallo an alle!}) an den Server.
        \item Empfangen Sie die Antwort und geben Sie sie auf der Konsole aus.
        \item Starten Sie \texttt{client.py} auf mehreren Rechnern gleichzeitig und tauschen Sie Nachrichten aus.
        \item Probieren Sie aus, ob Sie von verschiedenen Rechnern aus gleichzeitig Nachrichten senden und empfangen können. Was passiert, wenn mehrere Clients verbunden sind? Können Sie sich gegenseitig Nachrichten schicken?
        \item \textbf{Erweiterung:} Jeder Client gibt beim Start einen eigenen Namen ein. Beim Senden einer Nachricht wählen Sie aus, an welchen Namen (also an welchen anderen Client) die Nachricht geschickt werden soll. Die Nachricht wird dann nur an diesen Empfänger weitergeleitet.

    \end{enumerate}

    \textbf{Verwenden Sie folgende Grundstruktur:}

    \begin{lstlisting}[caption={client.py}]
    import socket

    client = socket.socket(socket.AF_INET, socket.SOCK_STREAM)
    client.connect(("88.198.193.239", 12345))  # IP-Adresse des Relay-Servers, Port 12345

    name = input("Dein Name: ")
    client.send(name.encode())  # Namen an Server senden

    def empfange():
        while True:
        antwort = client.recv(1024).decode()
        if not antwort:
            break
        print(antwort)

    import threading
    threading.Thread(target=empfange, daemon=True).start()

    while True:
        empfaenger = input("Empfänger-Name (oder 'exit'): ")
        if empfaenger.lower() == "exit":
        break
        nachricht = input("Nachricht: ")
        if nachricht.lower() == "exit":
        break
        # Format: EMPFAENGER:Nachricht
        client.send(f"{empfaenger}:{nachricht}".encode())

    client.close()
\end{lstlisting}
    \textbf{Hinweis:} Der Server muss die Namen der verbundenen Clients verwalten und Nachrichten entsprechend weiterleiten. So können Sie gezielt Nachrichten an einzelne Personen schicken.
    \begin{myanswer}
        \begin{lstlisting}[caption=\texttt{client.py}]
import socket
import threading

client = socket.socket(socket.AF_INET, socket.SOCK_STREAM)
client.connect(("88.198.193.239", 12345))  # IP-Adresse des Relay-Servers, Port 12345

name = input("Dein Name: ")
client.send(name.encode())  # Namen an Server senden

def empfange():
    while True:
        antwort = client.recv(1024).decode()
        if not antwort:
            break
        print(antwort)

threading.Thread(target=empfange, daemon=True).start()

while True:
    empfaenger = input("Empfänger-Name (oder 'exit'): ")
    if empfaenger.lower() == "exit":
        break
    nachricht = input("Nachricht: ")
    if nachricht.lower() == "exit":
        break
    # Format: EMPFAENGER:Nachricht
    client.send(f"{empfaenger}:{nachricht}".encode())

client.close()
        \end{lstlisting}
    \end{myanswer}
    \begin{myanswer}
        \begin{lstlisting}[caption=\texttt{server.py}]
import socket
import threading

server = socket.socket(socket.AF_INET, socket.SOCK_STREAM)
server.bind(("0.0.0.0", 12345))
server.listen()

clients = []
names = []

def handle_client(client):
    while True:
        try:
            nachricht = client.recv(1024).decode()
            if not nachricht:
                break
            empfaenger, nachricht = nachricht.split(":", 1)
            if empfaenger in names:
                index = names.index(empfaenger)
                clients[index].send(f"{names[clients.index(client)]}: {nachricht}".encode())
        except:
            break

    client.close()
    clients.remove(client)
    names.remove(names[clients.index(client)])

while True:
    client, addr = server.accept()
    clients.append(client)
    name = client.recv(1024).decode()
    names.append(name)
    print(f"{name} hat sich verbunden.")
    threading.Thread(target=handle_client, args=(client,), daemon=True).start()
server.close()
        \end{lstlisting}
    \end{myanswer}
\end{myexercise}

\begin{myexercise}
    [Reflexionsfragen]
    Notieren Sie nach Abschluss von Learning Task 3 (max. 5 Minuten) kurze Antworten in Ihr Lernjournal. Tauschen Sie sich danach mit einer Person aus und ergänzen Sie fehlende Punkte.
    \textbf{Fragen:}
    \begin{enumerate}
        \item Was ist der Unterschied zwischen einer lokalen IP-Adresse und einer öffentlichen IP-Adresse?
        \item Erklären Sie kurz, was NAT (Network Address Translation) ist und warum es in Heimnetzwerken verwendet wird.
        \item Warum ist Port-Forwarding notwendig, um Verbindungen von aussen zu einem Gerät im lokalen Netzwerk herzustellen?
        \item Was ist ein Relay-Server und welche Rolle spielt er in diesem Szenario?
        \item Wie funktioniert DNS (Domain Name System) und warum ist es wichtig für die Kommunikation im Internet?
        \item Kurzer Selbstcheck (Ja/Nein):
              \begin{itemize}
                  \item Ich kann den Code ohne Vorlage erneut schreiben.
                  \item Ich verstehe jede Zeile von \texttt{client.py}.
                  \item Ich weiss, wie ich die öffentliche IP-Adresse meines Netzwerks finde.
              \end{itemize}
        \item Reflektieren Sie: Welche Herausforderungen könnten auftreten, wenn Sie versuchen, über das Internet zu kommunizieren, anstatt nur im lokalen Netzwerk?
    \end{enumerate}
    \begin{myanswer}
        \begin{itemize}
            \item Lokale IP: Identifiziert ein Gerät im lokalen Netzwerk (z.B. 192.168.x.x); öffentliche IP: Identifiziert das gesamte Netzwerk im Internet.
            \item NAT übersetzt lokale IP-Adressen in eine öffentliche IP-Adresse, damit mehrere Geräte im Heimnetzwerk über eine einzige öffentliche IP-Adresse kommunizieren können.
            \item Port-Forwarding leitet eingehende Verbindungen von der öffentlichen IP-Adresse an die richtige lokale IP-Adresse und den entsprechenden Port im Heimnetzwerk weiter.
            \item Ein Relay-Server empfängt Nachrichten von Clients und leitet sie an die vorgesehenen Empfänger weiter, wodurch direkte Verbindungen zwischen Clients vermieden werden.
            \item DNS übersetzt Domain-Namen in IP-Adressen, damit Menschen sich leicht zu merkende Namen verwenden können, anstatt sich numerische IP-Adressen merken zu müssen.
            \item Selbstcheck: Code reproduzierbar? Jede Zeile verstanden? Öffentliche IP finden gelingt? Wenn Nein $\rightarrow$ gezielt nacharbeiten.
            \item Herausforderungen Internet: NAT und Port-Forwarding nötig; dynamische öffentliche IP; zusätzliche Latenz; Firewalls/Provider-Restriktionen; Sicherheit (Verschlüsselung, Authentifizierung); DNS-Auflösung.
        \end{itemize}
    \end{myanswer}
\end{myexercise}

\clearpage
\section{Lösugen der längeren Code-Aufgaben}
\begin{myanswer}[ex:chat-localhost]
    \lstinputlisting[caption={server\_localhost.py}]{Grundlagen_Info/07_Netzwerke/Code/0_localhost/server_localhost.py}
    \lstinputlisting[caption={client\_localhost.py}]{Grundlagen_Info/07_Netzwerke/Code/0_localhost/client_localhost.py}
\end{myanswer}
\newpage
\begin{myanswer}[ex:chat-localnet]
    \lstinputlisting[caption={\texttt{chat\_localnet\_client.py}}]{Grundlagen_Info/07_Netzwerke/Code/1_LAN/chat_localnet_client.py}
    \lstinputlisting[caption={\texttt{chat\_localnet\_server.py}}]{Grundlagen_Info/07_Netzwerke/Code/1_LAN/chat_localnet_server.py}
\end{myanswer}
\newpage
\begin{myanswer}[ex:endlos-chat]
    \lstinputlisting[caption={\texttt{chat\_localnet\_client\_loop.py}}]{Grundlagen_Info/07_Netzwerke/Code/1_LAN/chat_localnet_client_loop.py}
    \lstinputlisting[caption={\texttt{chat\_localnet\_server\_loop.py}}]{Grundlagen_Info/07_Netzwerke/Code/1_LAN/chat_localnet_server_loop.py}
\end{myanswer}